\documentclass[12pt]{article}

\usepackage{amsmath}
\usepackage{amssymb}
\usepackage{amsthm}
\usepackage{mathtools}
\usepackage{mathrsfs}
\usepackage{hyperref}
\usepackage{braket}
\usepackage{bm}
\usepackage{graphicx}
\usepackage[
    left = 1in,
    right = 1in,
    top = 1in,
    bottom = 1in
]{geometry}
\usepackage{titling}
\usepackage{parskip}
\usepackage{hyperref}
\usepackage{fancyhdr}

\pagestyle{fancy}
\fancyhf{}
\rhead{\small AC Multivariable Calculus}
\lhead{\small Final Project}
\cfoot{\thepage}
\renewcommand{\headrulewidth}{0.4pt}
\fancyfoot{}
\renewcommand{\footrulewidth}{0.4pt}
\rfoot{\small \thepage}
\lfoot{\small Charoenrattanaruk, Harpavat, Adams, and Meng}

\title{\textbf{Multivariable Calculus Final Project}}
\author{Parrie, Keshav, Diego, and Michael}
\date{}

\begin{document}

\maketitle
\thispagestyle{fancy}

\section*{Documentation}

\subsection*{Introduction \& Design}

For this project, our group chose a cone with radius $30"$ and height
$36"$ (or $2.5'$ and $3'$, respectively). The equation that we used was 
\begin{equation*}
    \frac{x^{2}}{900} + \frac{y^{2}}{900} = \frac{z^{2}}{1296},
\end{equation*}
where $x \in [-30, 30]$, $y \in [-30, 30]$, and $z \in [0, 36]$.

\newpage

\section*{Calculations}

\subsection*{Theoretical Precise Value}

To calculate the precise theoretic value, we will use the triple integral
\begin{equation*}
    \iiint\limits_{Q} f(x, y, z) \, dV.
\end{equation*}
To find the bounds, we will use cylindrical coordinates, as Cartesian
coordinates tend to end up messy in such calculations. To find the bounds
in cylindrical coordinates, we only need to solve for $r$, as $\theta \in [0,
2\pi]$ and $z \in [0, 36]$, as given by our domain restrictions above.

To find the bounds for $r$, we first convert the rectangular $x$ and $y$
to cylindrical coordiantes. Following the conversion, our cone can be 
represented by the equation
\begin{equation*}
    \frac{r^{2}}{900} = \frac{z^{2}}{1296}.
\end{equation*}
Solve for $r$:
\begin{equation*}
    r = \sqrt{ \frac{900z^{2}}{1296} } = \sqrt{ \frac{25z^{2}}{36} } = \frac{5z}{6}. 
\end{equation*}
Thus, the upper bound for $r$ is $\frac{5z}{6}$, and because the lower  
bound for $r$ is $0$, we know that the complete bounds for $r$ are $[0,
\frac{5z}{6}]$. Thus, the complete bounds, for all variables, are:
\begin{equation*}
    \theta \in \left[ 0, 2\pi \right], \ z \in \left[ 0, 36 \right], \ \text{and}
    \ r \in \left[ 0, \frac{5z}{6} \right].
\end{equation*}

Therefore, the volume integral for the cone is If
we plug in all the values for the triple integral, we end up with
\begin{equation*}
    \int_{0}^{2\pi} \int_{0}^{36} \int_{0}^{\frac{5z}{6}} r \, dr  \, dz  \, d\theta,
\end{equation*}
which we can solve as follows:
\begin{align*}
    &V = \int_{0}^{2\pi} \int_{0}^{36} \int_{0}^{\frac{5z}{6}} r \, dr  \, dz  \, d\theta
    = \int_{0}^{2\pi} \int_{0}^{36} \left[ \frac{r^{2}}{2} \right]_{0}^{\frac{5z}{6}} \, dz  \, d\theta
    = \int_{0}^{2\pi} \int_{0}^{36} \frac{25z^{2}}{72} \, dz  \, d\theta \\
    &= \int_{0}^{2\pi} \left[ \frac{25z^{3}}{216} \right]_{0}^{36} \, d\theta
    = 5040 \int_{0}^{2\pi} \, d\theta
    = 5040 \, (2\pi)
    = 10800 \, (2\pi)
    = 33929.2007.
\end{align*}

Thus, the theoretical precise volume of our cone is $\boxed{33929.2007 \
\text{cubic inches}}$ (or approxinately $19.635$ cubic feet).

\end{document}
